\begin{frame}[fragile]{xsl:if}
    \begin{block}{\textcolor{blue}{\texttt{<xsl:if \textcolor{orange}{test="chemin\_xpath"}>}}}
    \begin{itemize}
        \item Une condition qui applique les règles si \textcolor{orange}{\texttt{@test}} est vraie.
    \end{itemize}
    \begin{lstlisting}[numbers=none]
        <xsl:template match="body">
            <xsl:if test="lg[@type='sonnet']">
                <xsl:copy-of select="."/>
            </xsl:if>
        </xsl:template>
    \end{lstlisting}
        
    \end{block}
\end{frame}

\begin{frame}[fragile]{xsl:choose, xsl:when, xsl:otherwise}
    \begin{block}{\textcolor{blue}{\texttt{<xsl:choose>}}}
        \begin{itemize}
            \item Est utilisé avec au moins un élément \textcolor{blue}{\texttt{<xsl:when>}}, et en option un élément \textcolor{blue}{\texttt{<xsl:otherwise>}} pour exprimer plusieurs conditions / sélectionner une possibilité dans une liste de choix
        \end{itemize}
    
        \begin{lstlisting}[numbers=none]
            <xsl:template match="l">
                <xsl:choose>
                    <xsl:when test="placeName">
                       <xsl:text>LIEU</xsl:text>
                    </xsl:when>
                    <xsl:otherwise>
                        <xsl:text>PAS DE LIEU</xsl:text>
                    </xsl:otherwise>
                </xsl:choose>
            </xsl:template>
        \end{lstlisting}
        
    \end{block}
\end{frame}

\begin{frame}[fragile]{xsl:for-each, xsl:sort, xsl:for-each-group}
    \begin{block}{\textcolor{blue}{\texttt{<xsl:for-each>}}}
        \begin{itemize}
            \item Itère sur les noeuds sélectionnés par l'\textcolor{orange}{@select} et applique les règles à chacun des éléments du noeud.
        \end{itemize}

        \begin{lstlisting}
            <xsl:template match="l">
                <xsl:for-each select="placeName">
                    <xsl:text> TOTO </xsl:text>
                </xsl:for-each>
            </xsl:template>
        \end{lstlisting}
    \end{block}
\end{frame}

\begin{frame}[fragile]{xsl:for-each, xsl:sort, xsl:for-each-group}
    \begin{block}{\textcolor{blue}{\texttt{<xsl:sort>}}}
        \begin{itemize}
            \item Permet de changer l'ordre des noeuds avec les attributs suivants:
            \begin{itemize}
                \item \textcolor{orange}{@select}: objet du tri
                \item \textcolor{orange}{@data-type}: par défaut est alphabétique, mais on peut mettre "number"
                \item \textcolor{orange}{@order}: par défaut "ascending"
                \item \textcolor{orange}{@case-order}: indique si ce sont les minuscules ou les majuscules qui apparaîtront en premier.
            \end{itemize}
        \end{itemize}
        \begin{lstlisting}[numbers=none]
            <xsl:template match="element">
                <xsl:for-each select="souselements">
                    <xsl:sort select="tri" order='ascending|descending'/>
                    <xsl:copy-of select="."/>
                </xsl:for-each>
            </xsl:template>
        \end{lstlisting}
    \end{block}
\end{frame}

\begin{frame}[fragile]{xsl:for-each, xsl:sort, xsl:for-each-group}
    \begin{block}{\textcolor{blue}{\texttt{<xsl:for-each-group>}}}
        \begin{itemize}
            \item Itère sur les groupes de noeuds sélectionnés par \textcolor{orange}{\texttt{@select}} et applique la règle à chacun d'entre eux
            \item \textcolor{orange}{\texttt{@group-by}} permet de rassembler les noeuds sélectionnés en sous-groupes.
        \end{itemize}
        
    \end{block}
\end{frame}

\begin{frame}[fragile]{xsl:variable}
    \begin{block}{\textcolor{blue}{\texttt{<xsl:variable \textcolor{orange}{name="nom\_variable" select="'valeur\_variable'"}/>}}}
        \begin{itemize}
            \item Créer une variable dans la feuille de style.
            \item Utile pour la mise en forme par exemple.
        \end{itemize}
    \end{block}
    \begin{lstlisting}
        <xsl:variable name="header">
          <font color="#9acd32"> Title </font>
        </xsl:variable>

        <xsl:template match="/">
          <html>
          <body>
            <xsl:copy-of select="$header"/>
          </body>
          </html>
        </xsl:template>
    \end{lstlisting}
\end{frame}

